\documentclass[quick,con]{debate}
\motion{向孩子解释死亡}
\stance{反方}{更应该用故事}
\begin{document}
\debatetitle
\begin{thesisbox}[一句话立论]解释成功与否不看大人说了什么，而看孩子接住了什么，而在孩子最需要解释的那个年纪，他能接住的是故事。\end{thesisbox}
\section{定义与判准（标准表述）}
\begin{fields}
\field{故事}{以叙事、比喻、拟人为主要载体的解释，字面内容不要求与事实一一对应，但也不要求与事实冲突。
}
\field{现实}{以死亡的事实属性为解释的主要内容，语气可以温柔。这一点不跟对方争。
}
\field{孩子}{三到十二岁，重心三到八岁；这个年纪的关键事实是死亡概念还在成形。
}
\field{判准}{哪一种方式更能让孩子既正确地知道发生了什么，又承受得住它、并且愿意继续跟大人谈。重点在后半句：看孩子接住了什么，不是看大人说出了什么。
}
\field{对方提偏向定义或判准时的 30 秒口径}{对方把判准改成“孩子最终获得的概念是否正确”：按这条标准，最好的做法是给四岁孩子放解剖视频，评委不会同意；孩子问死亡从来不只是在要一个知识。对方把“故事”定义成“用不真实的说法代替事实”：那是定义杀，请先回答《獾的礼物》算不算故事。
}
\end{fields}
\section{我方论点}
\begin{longtable}{@{}L{\dimexpr0.062\linewidth-1.50\tabcolsep}L{\dimexpr0.155\linewidth-1.50\tabcolsep}L{\dimexpr0.490\linewidth-1.50\tabcolsep}L{\dimexpr0.293\linewidth-1.50\tabcolsep}@{}}
\toprule \thead{标签} & \thead{一句话} & \thead{核心数据 / 学理 / 案例} & \thead{出处}\\\midrule\endhead
孩子先懂故事 & 解释成功看孩子接住了什么；那个年纪能接住的形式是叙事 & 河南 114 名 5–6 岁儿童五子成分成熟率：不可逆性 67.8\%、因果性 60.7\%、非功能性 57.1\%、普遍性 55.3\%、死亡判定 53.6\% & Liu 与 Liu 2024，《Frontiers in Psychology》15\\
同上 & 叙事是更有效的传递形式 & 元分析：75 个以上样本、33000 多名参与者，故事比说明文更易理解、记得更牢（局限：测的是文本理解与记忆，不是死亡观，自己先说出来） & Mar 等 2021，《Psychon Bull Rev》28(3): 732–749\\
故事多给一样 & 孩子可以一边知道人死了不会回来，一边在故事里保留这段关系 & 墨西哥 61 名儿童：81.9\% 信亡者探望、78.6\% 信亡者进食，同时全部孩子都有生物学死亡理解；父母宗教与灵性思考对四个子成分均无显著影响 & Gutiérrez 等 2020，《Child Development》\\
同上 & 维系联结不是病理 & Klass、Silverman 与 Nickman 1996 持续联结理论 & 《Continuing Bonds》\\
案例 & 一本一句假话都没说的绘本 & 《獾的礼物》：獾死了，没复活也没去别处，朋友靠回忆活下去 & Varley 1984，Andersen Press\\
\bottomrule\end{longtable}
\section{对方论点 → 一句话反驳}
\begin{longtable}{@{}L{\dimexpr0.210\linewidth-1.33\tabcolsep}L{\dimexpr0.284\linewidth-1.33\tabcolsep}L{\dimexpr0.505\linewidth-1.33\tabcolsep}@{}}
\toprule \thead{对方会说} & \thead{我方一句话回} & \thead{对方追问时}\\\midrule\endhead
孩子问事实，你们答来世，答非所问 & 原文写的是家长“扩展”，不是“替换”，家长回答里 55.0\% 答的正是死因 & 同一份研究里全部孩子仍有生物学理解\\
孩子按字面理解比喻，产生新恐惧 & 那是敷衍的问题，不是故事的问题 & 请对方用《獾的礼物》举一次例；做不到，说明他反对的是敷衍不是叙事\\
训练实验证明讲现实更好 & 那节课教的是身体，场景是课堂，不是丧亲 & 60 人，准实验；而“越懂生物学越不怕死”是横断面相关，方向可能反过来\\
故事被拆穿，孩子不再信任大人 & 孩子并不把故事当成可被拆穿的事实陈述 & 英国研究的成人组没有共存式思维，只剩清楚的两套解释，他们没崩溃，他们长大了\\
解释的目的就是让孩子知道真相 & 那条判准评委不会接受 & 孩子问的还有“我还会不会被爱”“你会不会也走”，这两个问题事实答不了\\
\bottomrule\end{longtable}
\section{必问与必答}
\begin{points}{必问 （对方不答就点名，自由辩前两分钟问完）}
\item 请举一个《獾的礼物》级别的故事，告诉我它错在哪。
\item 五到六岁孩子里还有三分之一没掌握不可逆性，这三分之一怎么办？
\item 孩子问死亡的时候，他是不是也在问“你会不会有一天也不见了”？
\item 外婆今晚刚走，你方的第一句话是什么？第二句呢？（一辩稿抛出的那个问题的落地版）
\end{points}
\begin{points}{必答 （15 秒内正面回）}
\item “变成星星”字面上是真的吗？ → 不是。所以我方不拿它当范例，我方的范例是《獾的礼物》。
\item 孩子把故事当真了怎么办？ → 会纠正。任何教育都有后续，一次说不到位不等于方法错。
\item 你反对说“死了”这两个字吗？ → 不反对。我方反对的是说完就没了。
\item 亡灵节有 31.1\% 的孩子说害怕？ → 是，而且随年龄下降，四岁 55.0\%，六岁 15.0\%。怕本来就是死亡的一部分。
\item 你那份元分析测的不是死亡观吧？ → 对，立论里我方自己先说了，只用它说明叙事传递更有效。
\end{points}
\section{自由辩战场概要}
\begin{longtable}{@{}L{\dimexpr0.150\linewidth-1.33\tabcolsep}L{\dimexpr0.410\linewidth-1.33\tabcolsep}L{\dimexpr0.440\linewidth-1.33\tabcolsep}@{}}
\toprule \thead{战场} & \thead{开场问题} & \thead{收束语}\\\midrule\endhead
一 正方在反对什么 & 请举一个《獾的礼物》级别的故事，告诉我它错在哪 & 一整场下来，对方举的例子没有一个是故事\\
二 解释的那一半（主场） & 孩子是不是也在问“你会不会有一天也不见了”？ & 对方解释的是死亡这件事，我方解释的是孩子的那个夜晚\\
三 准确性被占用了吗 & 父母的宗教与灵性倾向对四个子成分都没有显著影响，占用在哪？ & 争的不是要不要说真话，是真话要用什么形状交到孩子手上\\
\bottomrule\end{longtable}
\section{如果……就……}
\begin{contingency}
\ifthen{对方改定义、改判准，或拿出没预料到的数据}{定义杀用《獾的礼物》做反例，判准按上面的 30 秒口径打；没见过的数据先问样本年龄上限与测的是什么变量。}
\ifthen{论点被打穿、对方抢走“温柔”、对方回避必问}{收缩到“解释的那一半”，回到判准；不争温柔，改口径；对方回避就点名、重复、不换题。}
\end{contingency}
\end{document}
